\section{Related Work}
\label{sec:related}

\subsection{Context layers}

MotherDuck reports 418 of 419 DABStep questions correct (99.8\%) with Gemini~3
Flash at roughly \$0.02 per question, and separately 96--98\% with a local
27B model, using ``Guides'' --- markdown context stored in the warehouse and
fetched through an MCP server~\cite{motherduck-guides}. Their headline claim is
the same as ours and measured the same way: context raises accuracy by 72
percentage points while cutting cost by roughly 55\%.

That is an independent corroboration of the direction of our finding, at a far
higher absolute score. Because a reader will otherwise stop at the difference
between 99.8\% and 56.6\%, we state plainly what that difference is made of.

\subsection{When the semantics get compiled}

The two approaches differ in \emph{when} a metric's definition is turned into
executable form.

\begin{center}
\small
\begin{tabular}{@{}p{0.28\columnwidth}p{0.28\columnwidth}p{0.28\columnwidth}@{}}
\toprule
 & \tabhead{Views / macros} & \tabhead{Declarative contract} \\
\midrule
Compiled & ahead of time, by a human & per query, by the model \\
Covers & anticipated metrics & any question in the domain \\
Cost & one macro per metric, maintained & one description per domain \\
Fails when & nobody anticipated the question & the model misapplies the rule \\
\bottomrule
\end{tabular}
\end{center}

MotherDuck's own guidance lists pre-computed views among its five steps, and a
separate post of theirs reports 93.2\% from ``a simple prompt with views and
macros''. \textbf{So a large share of the distance between our numbers and
theirs is work relocated from the agent to a human author, not reasoning
improved.}

We verified this directly rather than inferring it. Tasks 1500 and 1502, both
of which our contract arm answered incorrectly, are answered \emph{exactly} by
a five-line macro encoding the wildcard rule the contract states in prose:

\begin{lstlisting}[language=SQL]
CREATE MACRO wild(col, val) AS
  (col IS NULL OR len(col) = 0
   OR list_contains(col, val));
SELECT string_agg(ID::VARCHAR, ', '
                  ORDER BY ID)
FROM fees
WHERE wild(account_type, 'O') AND wild(aci, 'C');
\end{lstlisting}

The agent's job collapses from \emph{derive the rule} to \emph{call the macro}.

\subsection{Why the benchmark flatters pre-computation}

DABStep's 450 questions come from 26 templates, and 294 of 332 hard tasks in
our set are fee questions. Pre-built views are maximally effective when the
question space is enumerable in advance, which is exactly this benchmark's
construction. A macro contributes nothing to a question nobody wrote it for.

The declarative layer measured here is the layer that must cover the
\emph{unanticipated} question. \textbf{56.6\% is therefore a measurement of a
harder task, not a worse attempt at the same one.}

We want to be careful not to convert that into a claim we cannot support.
Maintaining one macro per metric per dialect is a real and growing burden, but
we have no data on where it crosses over against the cost of authoring domain
prose. \textbf{We do not claim the contract approach scales better.} We claim
only that the two degrade differently: a missing macro degrades to nothing,
while a contract degrades gracefully to a model reasoning from a description.

The generous reading is also the correct one. These are complementary layers,
not competitors --- macros optimise the head of the question distribution, a
contract covers the tail, and in a system with both, the contract is what tells
the agent which macro exists and when to call it.

\subsection{What is distinctive here}

\begin{center}
\small
\begin{tabular}{@{}p{0.29\columnwidth}p{0.21\columnwidth}p{0.29\columnwidth}@{}}
\toprule
 & \tabhead{Prior reports} & \tabhead{This work} \\
\midrule
Comparison & with vs without & four arms \\
Separates content from tooling & no & \textbf{yes}, via arm~D \\
Statistical test & two accuracies & paired McNemar, all pairs \\
Second model & no & \textbf{yes}, same direction \\
Artifact frozen pre-hoc & undisclosed & \textbf{yes}, digest-pinned \\
Transcripts released & repository & \textbf{3{,}208} \\
\bottomrule
\end{tabular}
\end{center}

The contribution is the decomposition, not the accuracy.

\subsection{Interpreting published leaderboard scores}

One observation from reconstructing golds bears on how any DABStep number
should be read. Across the 2{,}199 published submissions, organisations
submitting once achieve a median hard-split accuracy of 19.6\%; those
submitting 21 or more times achieve 54.2\%. Independent commentary on Guides
similarly warns against hardcoding answers and urges measuring generalisation
to unseen questions.

We draw no conclusion about any particular entrant --- a strong tuning gradient
is equally consistent with genuine iterative improvement --- and methodology is
not disclosed in enough detail for the 99.8\% run to rule anything in or out.
The point is structural: on a benchmark of 26 templates with a public
leaderboard, a single reported accuracy carries little information about
generalisation unless the submission count and the artifact's provenance are
stated. Ours are: one sweep per model, no tuning against the task set, and a
digest-pinned artifact frozen 35 commits before the first run.

\subsection{The deliberate non-response}

The obvious response to Section~\ref{sec:mechanism} is to make this library's
metric semantics executable --- a tool that runs a declared SQL expression
rather than describing it --- and re-run. That feature is worth building and is
on the roadmap for independent reasons.

It was deliberately \textbf{not} built before this measurement, because it
would be designed from knowledge of which tasks the contract arm lost. That is
fitting to the test set, and it would spend the strongest methodological asset
here: an artifact authored from the vendor manual alone and frozen before any
question was read. The order is measure, ship, re-measure. This paper is the
baseline that feature will be judged against.
